\documentclass[11pt]{article}
\usepackage[margin=1in]{geometry}
\usepackage[T1]{fontenc}
\usepackage{lmodern,amsmath,amssymb,amsthm,mathtools,hyperref}
\setlength{\emergencystretch}{3em}
\newtheorem{theorem}{Theorem}[section]
\newtheorem{lemma}[theorem]{Lemma}
\newtheorem{corollary}[theorem]{Corollary}
\newtheorem{proposition}[theorem]{Proposition}
\newcommand{\tr}{\operatorname{tr}}
\title{Sharp collision-entropy bounds and exact rounding\\for pairwise-independent sign variables}
\author{Joshua Mark\thanks{ORCID: \href{https://orcid.org/0009-0003-5905-2382}{0009-0003-5905-2382}. AI tools assisted with derivation, independent review, computation and exposition.}}
\date{20 November 2026 --- research draft; priority not established}
\begin{document}\maketitle
\begin{abstract}
We prove a residue-sensitive collision-probability bound for unbiased
pairwise-independent sign variables, with a complete equality
classification by Hadamard completion in at most three columns.
We also derive a fixed-dimension sufficient condition for a probability
distribution on sign vectors to contain an entire Hadamard matrix.
The condition combines collision probability with an explicit error
in its first and second moments. It yields positive Shannon-entropy
windows for exact extraction, and a larger existence window using
a classical seven-row completion theorem. We give a sharp collision
benchmark at forbidden orders and explicit obstructions to diagonal
balancing, natural operator capacities, convex exchangeable averaging,
and entrywise modular lifting. These results distinguish genuine
rounding certificates from relaxations that erase the sign constraints.
No construction attaining the sufficient window in every admissible
dimension is proved; the full Hadamard conjecture remains open.
Novelty has not been established by a comprehensive prior-art review.
\end{abstract}
\section{A collision-entropy bound in every residue class}
\label{sec:residue-bound}
Let $d\ge2$ unbiased pairwise-independent sign variables have joint
law $w$ on $\{\pm1\}^d$. Set
\[
 n=d+1,\qquad N=4\left\lceil\frac n4\right\rceil,
 \qquad C(w)=\sum_xw_x^2.
\]
The integer $N$ is the smallest multiple of four at least $d+1$.
The standard orthogonal-frame argument gives the atom bound
$w_x\le1/n$; the entropy connection and its exact $N=n$ case are
background~\cite{GP}. The residue-sensitive statement below is a
derivation in this research pass; bibliographic priority is unestablished.

\begin{theorem}[Residue-sensitive collision bound and equality]
\label{thm:residue-bound}
For every such law,
\[
 C(w)\le\frac1N,\qquad
 H_2(w):=-\log C(w)\ge\log N,\qquad H(w)\ge\log N.
\]
If $C(w)=1/N$, the law is uniform on exactly $N$ atoms and their
augmented sign matrix admits completion by $N-n\le3$ columns to a
Hadamard matrix of order $N$. Conversely any order-$N$ Hadamard matrix
gives a law attaining both entropy bounds. Consequently attainment
of either entropy lower bound is equivalent to Hadamard existence
at order $N$.
\end{theorem}

\begin{proof}
Write $v_x=(1,x)\in\{\pm1\}^n$. Unbiased pairwise independence says
\begin{equation}\label{eq:residue-frame}
 \sum_xw_xv_xv_x^T=I_n.
\end{equation}
Thus $\sum_yw_y(v_x\cdot v_y)^2=n$ for each $x$, and
$\sum_{x,y}w_xw_y(v_x\cdot v_y)^2=n$.
In particular $w_x\le1/n$.

\emph{Case $n\equiv0\pmod4$.}
The atom bound gives $C\le1/n$. Equality forces every positive weight
to equal $1/n$. The square augmented matrix has orthogonal columns
by \eqref{eq:residue-frame}, and is Hadamard.

\emph{Case $n\equiv3\pmod4$.}
All dot products are odd. Distinct augmented atoms cannot be negatives
of one another, and their squared dot products are at least one.
Consequently
\[
 n\ge n^2C+(1-C),\qquad C\le\frac1{n+1}.
\]
Equality forces every distinct pair to have dot product $\pm1$.
The pointwise identity then reads
$n=n^2w_x+1-w_x$, forcing $w_x=1/(n+1)$.
Let $V$ have these $N=n+1$ augmented rows. Then $V^TV=NI_n$;
$K=NI_N-VV^T$ is positive semidefinite of rank one with diagonal
one. Hence $K=ss^T$ for a sign vector $s$. Appending $s$ gives
$[V\ s][V\ s]^T=NI_N$.

\emph{Case $n\equiv2\pmod4$.}
Here $n\ge6$. Divide the atoms by their parity
$\epsilon_x=\prod_i x_i$; let their total masses be $a,1-a$.
Two atoms in the same class have dot product $2\pmod4$, while
cross-class dot products are $0\pmod4$. Therefore
\[
 n\ge n^2C+4\big(a^2+(1-a)^2-C\big)
       \ge(n^2-4)C+2,
\]
giving $C\le1/(n+2)$. Equality forces $a=1/2$, within-class
distinct dot products of magnitude two, and cross-class orthogonality.
For a fixed atom the pointwise identity is then
$n=n^2w_x+4(1/2-w_x)$, so $w_x=1/(n+2)$.
For the resulting $N\times n$ matrix $V$, the positive semidefinite
residual $K=NI_N-VV^T$ has rank two and diagonal two. Cross-class
blocks vanish, so each of the two nonempty diagonal blocks has rank
one. Write them $2ss^T$ and $2tt^T$, with $s,t$ sign vectors.
Append rows $(s_i,s_i)$ on the first class and $(t_i,-t_i)$ on the
second. These two added columns have row Gram matrix $K$, completing
$V$ to a Hadamard matrix.

\emph{Case $n\equiv1\pmod4$.}
Here $n\ge5$. For any pair of atoms set
$q_{xy}=\epsilon_x\epsilon_y(v_x\cdot v_y)$.
Their Hamming distance shows $q_{xy}\equiv1\pmod4$. Thus
\[
 q_{xy}^2+2q_{xy}-3\ge0,
\]
with equality precisely at $q_{xy}\in\{-3,1\}$.
Put $z=\sum_xw_x\epsilon_xv_x$. Since the $n$ coordinate functions
of $v$ are orthonormal in $L^2(w)$, Bessel inequality gives
$\|z\|^2\le\mathbb E\epsilon^2=1$. Equivalently, this follows by
expanding $\mathbb E(\epsilon-z\cdot v)^2=1-\|z\|^2$.
Summing the polynomial over pairs and retaining its diagonal terms gives
\[
 (n-1)(n+3)C\le n+2\|z\|^2-3\le n-1,
\]
so $C\le1/(n+3)$.

At equality, $\epsilon_x=z\cdot v_x$ on the positive-weight support,
and distinct pairs have $q_{xy}\in\{-3,1\}$. The polynomial's
pointwise expectation equals
$n+2\epsilon_xv_x\cdot z-3=n-1$; its only nonzero contribution
comes from $y=x$. Hence again $w_x=1/(n+3)$.
Switch row $v_x$ to $v'_x=\epsilon_xv_x$ and form $V'$.
Its residual $K=NI_N-V'V'^T$ is positive semidefinite of rank three,
has diagonal three, and off-diagonal entries in $\{3,-1\}$.
Represent it as a Gram matrix of length-$\sqrt3$ vectors. Inner
product three identifies equal vectors. Distinct types have Gram
matrix $4I_t-J_t$, which is positive semidefinite only if $t\le4$.
Assign the types distinct rows of
\[
 T=\begin{pmatrix}1&1&1\\1&-1&-1\\-1&1&-1\\-1&-1&1\end{pmatrix}.
\]
The resulting $N\times3$ sign matrix $R$ has $RR^T=K$.
Then $[V'\ R]$ is a Hadamard matrix. This is an explicit completion;
no external small-rank classification is needed.

For the converse normalize the first column of an order-$N$ Hadamard
matrix to all positive and keep $n$ columns including that column.
The $N$ resulting rows are distinct: full rows differ in $N/2$
coordinates, whereas at most three columns were deleted. For $N\ge8$
this leaves a difference; at $N=4$, $n\ge3$ and at most one column
was deleted. The uniform law on the remaining $d$ coordinates is
unbiased and pairwise independent, with $C=1/N$ and $H=\log N$.
Finally $H\ge-\log C$ by Jensen inequality, so equality in the Shannon
lower bound forces $C=1/N$ as well.
\end{proof}

For example, four, five or six unbiased pairwise-independent sign bits
each have Shannon entropy at least $\log8$. Eight, nine or ten such
bits each have entropy at least $\log12$, attained by truncating an
order-$12$ Hadamard matrix. These statements concern joint entropy,
not min-entropy; a single atom can still saturate $1/(d+1)$.
For $d=1$ the formula would be false, which is why $d\ge2$ is explicit.

\paragraph{The unresolved step.}
This proves a universal bound and a complete equality classification.
It does not prove that an equality law exists in every admissible
dimension. Such attainment for all $d$ is equivalent to the full
Hadamard conjecture. The parity and Bessel argument gives a precise
information-theoretic target; it does not by itself reach that target.
\section{A quantitative discrete rounding bridge}
\label{sec:discrete-rounding}
This section gives a sufficient fixed-dimension certificate, rather than
an existence proof. It combines the weighted clique inequality of
Motzkin--Straus with an exact second-moment identity. The atom bound
and its equality characterization are established background; see
Gavinsky--Pudl\'ak, \emph{On the joint entropy of $d$-wise-independent
variables}, Theorem 3.2 and Section 4
(\url{https://users.math.cas.cz/~pudlak/pairwise.pdf}).
The present quantitative combination is a derivation, not a priority claim.

Let $n\ge4$ be even, $\Omega=\{-1,1\}^{n-1}$, and let $w$ be any
probability measure on $\Omega$. Write
\[
 v_x=(1,x_1,\ldots,x_{n-1})\in\{-1,1\}^n,\quad
 M=\sum_xw_xv_xv_x^T,\quad C=\sum_xw_x^2,\quad
 E=\|M-I_n\|_F^2.
\]
Here $C$ is the probability that two independent draws coincide.
The diagonal of $M$ is always one. Therefore $E$ is exactly twice
the sum of the squared biases and squared pair correlations:
\[
 E=2\sum_i(\mathbb E x_i)^2+
   2\sum_{i<j}(\mathbb E x_ix_j)^2.
\]
Balanced pairwise independence is equivalent to $E=0$.

\begin{theorem}[Collision--moment rounding certificate]
\label{thm:discrete-rounding}
If
\begin{equation}\label{eq:discrete-certificate}
 (n^2-4)C>n+E-\frac4{n-1},
\end{equation}
then the support of $w$ contains $n$ atoms whose augmented vectors are
mutually orthogonal. These vectors give a Hadamard matrix of order $n$.
They can be extracted by repeated weight merging, without adding atoms.
More generally, suppose every distinct nonorthogonal pair in the support
has $|v_x\cdot v_y|\ge a$, where $0<a<n$. Then the sufficient condition is
\[
 (n^2-a^2)C>n+E-\frac{a^2}{n-1}.
\]
If the support lies in a single parity class and $4\mid n$, one may
use $a=4$ (for $n>4$), improving the constant from $4$ to $16$.
\end{theorem}
\begin{proof}
Join distinct support atoms when their augmented vectors are orthogonal,
and let $A=\sum_{x,y:\,v_x\perp v_y}w_xw_y$ be the ordered edge weight.
Every clique has size at most $n$. Exact expansion gives
\[
 \sum_{x,y}w_xw_y(v_x\cdot v_y)^2=\tr(M^2)=n+E.
\]
The diagonal contributes $n^2C$. Distinct nonedges contribute at least
$a^2(1-C-A)$, so
\[
 A\ge1-\frac{n+E}{a^2}+\frac{n^2-a^2}{a^2}C.
\]
Under the displayed hypothesis this exceeds $1-1/(n-1)$.
For completeness, the weighted clique inequality follows constructively:
if positive-weight vertices $x,y$ are not adjacent, transfer their combined
weight to whichever has larger weighted neighbor sum. The ordered edge
weight does not decrease, and one positive-weight vertex disappears.
Iterating yields a clique of size $r$, with edge weight
$1-\sum_{i=1}^r u_i^2\le1-1/r$. Hence the resulting clique must have
$r\ge n$, and orthogonality forces $r=n$. This is the classical
Motzkin--Straus merging argument; its original source is
\emph{Maxima for graphs and a new proof of a theorem of Tur\'an}
(1965), \url{https://doi.org/10.4153/CJM-1965-053-6}.
Finally, dot products are even when $n$ is even. Within one parity class,
Hamming distances are even, so dot products are multiples of four if
$4\mid n$.
\end{proof}

\begin{corollary}[Explicit collision gap]
\label{cor:discrete-gap}
An exactly balanced pairwise independent measure on $n-1$ bits with
\[
 C>\frac1n-\frac4{n(n-1)(n^2-4)}
\]
contains an order-$n$ Hadamard matrix in its support. Consequently,
if no such matrix exists, every such measure satisfies the opposite
nonstrict inequality. With single-parity support, replace both occurrences
of $4$ in the correction by $16$, for $n>4$.
\end{corollary}
The standard atom estimate $w_x\le1/n$ implies $C\le1/n$;
the corollary gives an explicit one-sided neighborhood below equality
in which exact extraction already works. The correction is of order
$n^{-4}$. We do not claim optimality. In particular, the uniform cube
has $E=0$ but $C=2^{-(n-1)}$, far outside this neighborhood.
Thus positivity of every moment matrix, or exact vanishing of all
nonconstant Walsh coefficients of degrees one and two, does not itself
supply this concentration certificate. The unresolved step is producing
sufficiently concentrated measures for every $4\mid n$.

\subsection{A sharp obstruction to entrywise finite-field lifting}
A modular orthogonality route needs an additional global repair mechanism.
Its elementary seed construction below is standard in modular-Hadamard
literature; see Kuperberg, \emph{Hadamard matrices modulo $p$ and small
modular Hadamard matrices}, Lemma 2.3,
\url{https://arxiv.org/abs/1409.0148}.
\begin{proposition}[Nonsingular modular seeds need not lift]
\label{prop:discrete-modular}
Let $m\ge3$, $n>4$, and $n\equiv4\pmod m$, and put $B=J_n-2I_n$.
Then $BB^T\equiv nI_n\pmod m$, but $B$ is not a real Hadamard matrix.
Moreover there is no real Hadamard matrix entrywise congruent to $B$
modulo $m$. For every odd prime $p$, the choices $n=4+4pt$ with $t\ge1$
give infinitely many admissible orders with an invertible modular seed
that has no such sign lift.
\end{proposition}
\begin{proof}
Direct multiplication gives $BB^T=4I_n+(n-4)J_n$.
Its offdiagonal entries are the nonzero integer $n-4$, and are zero
modulo $m$. Since $m\ge3$, the two signs have distinct residues modulo
$m$, so the only possible entrywise sign lift of $B$ is $B$ itself.
For $n=4+4pt$, $n\equiv4\ne0\pmod p$; thus modular orthogonality
implies invertibility over $\mathbb F_p$.
\end{proof}
This refutes entrywise lifting, not unrestricted existence of another
Hadamard matrix at the same order. It remains a counterexample after
requiring admissible order, full modular rank, and exact sign entries.
At $n=16,p=3$ the seed has every offdiagonal Gram entry equal to $12$.
The unique sign lift is therefore already obstructed modulo $9$.
More generally the obstruction occurs immediately modulo $p^2$ whenever
$p$ does not divide $4t$.

\subsection{A sharp collision obstruction in the forbidden parity}
\begin{proposition}[Sharp weighted collision bound]
\label{prop:discrete-parity-collision}
Let $n\ge6$ with $n\equiv2\pmod4$. Every balanced pairwise independent
measure on $n-1$ bits satisfies
\[
 C\le\frac1{n+2}.
\]
Whenever an order-$(n+2)$ Hadamard matrix exists, this bound is attained.
In particular $C\le1/8$ for five bits is sharp, whereas the general
atom estimate alone gives only $C\le1/6$.
\end{proposition}
\begin{proof}
Three sign vectors of length $n\equiv2\pmod4$ cannot be mutually
orthogonal: normalize the first to all ones; the other two balanced
vectors would have to share exactly $n/4$ positive positions.
Thus the orthogonality graph is triangle-free and the merging argument
gives $A\le1/2$. Substituting $E=0,a=2$ into the energy bound yields
\[
 n\ge n^2C+4(1-C-A)\ge(n^2-4)C+2,
\]
which rearranges to the assertion. To attain it, normalize the first
column of an order-$(n+2)$ Hadamard matrix and delete two other columns.
The $n+2$ resulting length-$n$ rows are distinct: full rows differ in
$(n+2)/2>2$ coordinates. Their uniform measure has balanced pairwise
independent remaining $n-1$ coordinates and collision probability
$1/(n+2)$.
\end{proof}
This sharp example calibrates the relaxation: exact degree-two moments
can be achieved with only two more equal-weight atoms than the
rank lower bound, even when the requested order is impossible. The
near-equality certificate above uses a substantially smaller collision
deficit. It supplies a verifiable sufficient bridge, but not a theorem
that the bridge is attainable in every admissible dimension.


\begin{theorem}[Equality rigidity and two-column reconstruction]
\label{thm:discrete-parity-equality}
For $n\ge6$, $n\equiv2\pmod4$, a balanced pairwise independent measure
on $n-1$ bits attains $C=1/(n+2)$ if and only if it is uniform on
$n+2$ distinct atoms whose augmented matrix can be completed by two
sign columns to an order-$(n+2)$ Hadamard matrix.
In particular the collision bound is attained if and only if an
order-$(n+2)$ Hadamard matrix exists.
\end{theorem}
\begin{proof}
Assume equality. The preceding energy bound and $A\le1/2$ must
both be equalities. Hence $A=1/2$, and every distinct nonorthogonal
support pair has dot product of magnitude two.
Partition the augmented vectors by the parity of their number of
negative entries. When $n\equiv2\pmod4$, same-parity dot products
are congruent to two modulo four, and opposite-parity dot products
are divisible by four. Thus same-parity pairs cannot be orthogonal,
and every opposite-parity pair must be orthogonal: otherwise its
nonzero magnitude would exceed two. The orthogonality graph is
complete bipartite. Its two classes have masses $t$ and $1-t$;
$A=2t(1-t)=1/2$ forces both masses to be $1/2$.
For each support atom $x$, the pointwise identity $v_x^TMv_x=n$
now reads
\[
 n=n^2w_x+4(1/2-w_x),
\]
so $w_x=(n-2)/(n^2-4)=1/(n+2)$. The support therefore has
$m=n+2$ atoms. Form the $m\times n$ matrix $V$ with these rows.
Exact moments imply $V^TV=mI_n$.

Here is a self-contained completion argument. The residual matrix
$D=mI_m-VV^T$ is positive semidefinite of rank two and has diagonal
two. The opposite-parity orthogonality makes $D$ block diagonal on
the two nonempty parity classes. Each diagonal block has positive
rank; their ranks sum to two, so each has rank one. A positive
semidefinite rank-one matrix with diagonal two has form $2ss^T$
for a sign vector $s$. Write the two blocks as $2ss^T$ and $2tt^T$.
Append to $V$ the two columns whose rows are $(s_i,s_i)$ on the
first class and $(t_j,-t_j)$ on the second. Their Gram matrix is
exactly $D$, so the resulting square sign matrix has row Gram
$mI_m$ and is Hadamard. Conversely, deleting two nonconstant columns
from any normalized order-$m$ Hadamard matrix gives the required
uniform measure, as proved above.
\end{proof}
The equality statement is a rigidity theorem and an exact reconstruction
algorithm. It does not establish that the collision maximum reaches
this value when existence at the adjacent admissible order is unknown.
\section{Entropy concentration: a sufficient bridge with a measurable gap}
\label{sec:entropy-bridge}
We use natural logarithms. Let $\mu$ be a balanced pairwise-independent
probability measure on $n-1$ sign bits. Put
$H(\mu)=-\sum_x\mu_x\log\mu_x$ and $C(\mu)=\sum_x\mu_x^2$.
The background equality $H(\mu)=\log n$ if and only if its support
is an equally weighted Hadamard design is established by
Gavinsky--Pudl\'ak, Theorem 4.1~\cite{GP}. The quantitative implication
below follows from the collision certificate of
Theorem~\ref{thm:discrete-rounding}; we do not assert priority.

\begin{corollary}[A positive entropy window suffices]\label{cor:entropy-window}
For even $n\ge4$, define
\[
 \varepsilon_n=-\log\left(1-\frac4{(n-1)(n^2-4)}\right)>0.
\]
If $H(\mu)<\log n+\varepsilon_n$, the support of $\mu$ contains
$n$ mutually orthogonal augmented sign vectors. In particular it
contains an order-$n$ Hadamard matrix. Here
$\varepsilon_n\sim4/n^3$.
\end{corollary}
\begin{proof}
Jensen inequality gives
$\log C=\log\mathbb E_{x\sim\mu}\mu_x
\ge\mathbb E_{x\sim\mu}\log\mu_x=-H(\mu)$.
The assumed entropy bound therefore implies
\[
 C>\frac1n-\frac4{n(n-1)(n^2-4)}.
\]
Apply Corollary~\ref{cor:discrete-gap}.
\end{proof}
This is a fixed-dimension extraction criterion with a nonzero window;
it is not a construction of measures in that window. It changes an
exact entropy target to an explicit approximation target, but attaining
that target for all admissible orders is still unproved.

\subsection{Why optimal min-entropy is insufficient}
The entropy in this criterion cannot be replaced by min-entropy
$H_{\min}(\mu)=-\log\max_x\mu_x$.
\begin{proposition}[One saturated atom in every even dimension]
\label{prop:minentropy-counter}
For every even $n\ge4$ there is a balanced pairwise-independent measure
on $n-1$ bits with $H_{\min}=\log n$. At $n=6$ this holds although
no Hadamard matrix of order $6$ exists.
\end{proposition}
\begin{proof}
Put mass $1/n$ on the all-positive vector and spread the remaining
mass $(n-1)/n$ uniformly on the slice
\[
 \left\{x\in\{\pm1\}^{n-1}:\sum_i x_i=-1\right\}.
\]
Within the slice, permutation symmetry gives
$\mathbb E x_i=-1/(n-1)$ and
$\mathbb E x_ix_j=-1/(n-1)$ for distinct $i,j$, the latter by
expanding $(\sum x_i)^2=1$. Mixing in the positive atom makes both
moments zero. The slice has $\binom{n-1}{(n-2)/2}\ge n-1$ elements,
so no other atom has mass exceeding $1/n$.
At $n=6$ the ten slice atoms each have mass $1/12$, and the positive
atom has mass $1/6$. Nevertheless a real Hadamard matrix of order
$6$ is impossible: three orthogonal sign rows, after normalizing
the first to all positive, would force their two positive supports
to intersect in $6/4$ positions.
\end{proof}
Thus an information-theoretic bridge must control collision or Shannon
entropy, not just the mass of the heaviest outcome. The construction
and the implication are derived here; the general atom bound
$\mu_x\le1/n$ is background~\cite{GP}.

\subsection{A larger self-contained window when new rows are allowed}
If only existence is required, an orthogonal clique of $n-1$ rows
already suffices. Indeed, for its row matrix $X$,
$P=I-X^TX/n$ is a rank-one orthogonal projection with every diagonal
entry equal to $1/n$. Writing $P=zz^T$ with $\|z\|=1$ forces
$|z_j|=n^{-1/2}$; the row $\sqrt n z^T$ is the final sign row.
Using the merging argument with target clique size $n-1$ gives the
self-contained sufficient condition
\[
 (n^2-4)C>n+E-\frac4{n-2}.
\]
At $E=0$, this is
\[
 C>\frac1n-\frac8{n(n-2)(n^2-4)}.
\]
It is also implied by
$H(w)<\log n-\log(1-8/((n-2)(n^2-4)))$.
The extra row need not lie in the original support. This entropy
window is asymptotic to $8/n^3$ and needs no external classification.

\subsection{A stronger existence window using a classical completion theorem}
Verheiden's 1978 thesis, Theorem 4.5 on printed pages 22--23,
states that every $(n-r)\times n$ partial Hadamard matrix, with
$4\mid n$ and $r\le7$, completes to a full Hadamard
matrix~\cite{Verheiden}. We inspected the original scanned statement
and proof discussion. Its small-rank classification uses computation;
we have not independently reproduced that classification. The next
corollary explicitly uses this external mathematical input.

\begin{corollary}[Completion-assisted rounding]\label{cor:completion-rounding}
Let $4\mid n$, $n\ge12$, and let $w$ be any probability measure on
$n-1$ sign bits, with $C$ and moment defect $E$ as in
Theorem~\ref{thm:discrete-rounding}. If
\[
 (n^2-4)C>n+E-\frac4{n-8},
\]
then a Hadamard matrix of order $n$ exists. If $E=0$, it suffices that
\[
 C>\frac1n-\frac{32}{n(n-8)(n^2-4)}.
\]
Equivalently a sufficient Shannon-entropy window is
\[
 H(w)<\log n-\log\left(1-\frac{32}{(n-8)(n^2-4)}\right).
\]
The entropy conclusion here assumes $E=0$.
\end{corollary}
\begin{proof}
The ordered orthogonality edge weight in the proof of
Theorem~\ref{thm:discrete-rounding} exceeds $1-1/(n-8)$.
Its deterministic weight-merging argument therefore extracts a clique
of at least $n-7$ augmented sign vectors. Select $n-7$ of them and
apply the stated completion theorem. Algebra gives the collision
threshold; Jensen inequality gives the entropy version.
\end{proof}
Unlike the self-contained certificate, this conclusion allows the
last seven rows to be outside the original support. The correction
is asymptotic to $32/n^4$ in collision probability and $32/n^3$ in
Shannon entropy. It still supplies no measure attaining the window.
Weight merging is an extraction step; it does not preserve the
pairwise-independence constraints and cannot be iterated as a
moment-preserving construction without a separate argument.
\section{A lattice-theoretic lead and a sharp small-defect boundary}
\label{sec:completion-lead}
The classical completion theorem suggests studying the orthogonal
complement of a partial Hadamard matrix through arithmetic quadratic
forms. It is not an automatic solution: rational orthogonality,
integrality and sign entries are different requirements.

Let $X$ have $n-r$ orthogonal sign rows of length $n$, so
$XX^T=nI$. A real orthogonal completion $Y$ satisfies
\[
 Y^TY=nI-X^TX.
\]
The right side has rank $r$, diagonal entries $r$, integer entries,
and rigid congruences inherited from signs. Completing $X$ to a
Hadamard matrix asks for a factorization of this specific Gram matrix
by an $r\times n$ \emph{sign} matrix. A general real, rational, or
integer Gram factorization is a weaker problem.

\begin{proposition}[Eight missing rows can already be fatal]
\label{prop:eight-defect}
Let $K$ be a Hadamard matrix of order $m$ and let $s>1$ be odd.
Repeating every column $s$ times gives an $m\times ms$ partial
Hadamard matrix with no additional orthogonal sign row.
In particular there is an unextendible $4\times12$ example, with
eight missing rows, although Hadamard matrices of order $12$ exist.
\end{proposition}
\begin{proof}
The repeated matrix is $X=K\otimes(1,\ldots,1)$, and
$XX^T=msI_m$. For a putative new sign row, let $b_j$ be its sum
on the $j$th length-$s$ block. Orthogonality says $Kb=0$, so $b=0$
by invertibility of $K$. Each $b_j$ is a sum of an odd number of
signs and cannot be zero. Take $m=4,s=3$ for the example.
\end{proof}
This example is classical: the $4\times12$ pattern is displayed
in Verheiden's Section 4~\cite{Verheiden}. Our proof makes its block
obstruction explicit. Combined with his seven-row theorem, it shows
that universal completion based only on the number of missing rows
cannot be pushed from seven to eight.

The thesis also uses arithmetic results of Kneser and others to
restrict denominators of rational Gram completions. This is a
concrete route for further research: classify the residual Gram
forms and identify additional structure permitting sign factorizations
at rank eight and beyond. An $E_8$ explanation or an all-rank
completion theorem is not established here. In particular, the
unextendible example shows that any useful theorem must select or
modify the partial rows, rather than complete every possible choice.
\section{Why standard balancing and capacity do not force flatness}
\label{sec:capacity-obstruction}

Operator scaling connects invariant theory, noncommutative rank and
capacity; primary sources include Garg--Gurvits--Oliveira--Wigderson,
\emph{Operator scaling: theory and applications},
\url{https://arxiv.org/abs/1511.03730}, and Franks,
\emph{Operator scaling with specified marginals},
\url{https://arxiv.org/abs/1801.01412}. These theorems suggest a natural
test: can balancing the entry magnitudes of an orthogonal matrix force
a flat representative? The following elementary obstructions identify
precisely what that test loses. The arguments below are independently
derived; no external implementation is used.

For $U\in O(n)$ write $B_{ij}=u_{ij}^2$. Then $B$ is doubly stochastic:
it is nonnegative and every row and column sums to one. A matrix of
this form is called orthostochastic. Flatness asks for the very special
orthostochastic matrix $J_n/n$, not merely for prescribed marginals.
The distance-minimization formulation itself is established literature;
see Arasu--Mohan, \emph{Entropy of orthogonal matrices and minimum
distance orthostochastic matrices from the uniform van der Waerden
matrices}, Discrete Optimization 31 (2019), 115--144,
\url{https://doi.org/10.1016/j.disopt.2018.10.001}.

\begin{proposition}[Rigidity of positive diagonal balancing]
Let $U\in O(n)$ and let $D,E$ be positive diagonal matrices.
If $DUE\in O(n)$, then $DUE=U$ entry by entry.
If the bipartite support graph of $U$ is connected, necessarily
$D=cI_n$ and $E=c^{-1}I_n$ for some $c>0$.
\end{proposition}
\begin{proof}
Put $a_i=D_{ii}^2$, $b_j=E_{jj}^2$, and
$B'_{ij}=a_iB_{ij}b_j$. Both $B,B'$ are doubly stochastic and have
the same support. For equally normalized nonnegative arrays, relative
entropy is nonnegative, with equality exactly when the arrays agree.
Here direct summation using the marginals gives
\[
 \sum_{ij:B_{ij}>0}B'_{ij}\log\frac{B'_{ij}}{B_{ij}}
 =\sum_i\log a_i+\sum_j\log b_j,
\]
whereas the reverse relative entropy is the negative of that expression.
Both are nonnegative, so both vanish and $B'=B$.
Positive diagonal multipliers preserve the signs of every nonzero
entry, proving $DUE=U$. On each support edge $a_ib_j=1$;
connectedness propagates a common value of $a_i$ and its reciprocal
among the $b_j$, giving the final statement.
\end{proof}

This does not claim that general left-right matrix scaling is powerless.
It says that positive \emph{diagonal} scaling constrained to preserve
orthogonality cannot change the matrix at all.

\begin{proposition}[The natural classical and operator capacities are blind]
Let $B$ be any doubly stochastic matrix. Define
\[
 p_B(x)=\prod_{i=1}^n\left(\sum_j B_{ij}x_j\right),\qquad
 \operatorname{cap}(p_B)=\inf_{x_j>0}\frac{p_B(x)}{\prod_jx_j}.
\]
Also define the completely positive map
\[
 T_B(X)=\operatorname{diag}\bigl(B\operatorname{diag}(X)\bigr),
 \qquad
 \operatorname{cap}(T_B)=\inf_{X\succ0}
       \frac{\det T_B(X)}{\det X}.
\]
Then both capacities equal $1$. Moreover
$T_B(I)=T_B^*(I)=I$. Thus every orthogonal $U$, flat or not,
already passes these balancing and capacity tests for $B=U\circ U$.
\end{proposition}
\begin{proof}
Weighted arithmetic--geometric mean, followed by the column sums, gives
\[
 \prod_i\sum_jB_{ij}x_j
 \ge\prod_{ij}x_j^{B_{ij}}=\prod_jx_j.
\]
Equality is achieved at $x=(1,\ldots,1)$, proving the polynomial claim.
The Kraus matrices $K_{ij}=\sqrt{B_{ij}}E_{ij}$ represent $T_B$, so it
is completely positive. For $X\succ0$, the preceding inequality and
the determinant bound by the product of diagonal entries yield
\[
 \det T_B(X)\ge\prod_jX_{jj}\ge\det X.
\]
Taking $X=I$ proves the operator-capacity assertion. The row and column
sums give the two marginal identities directly.
\end{proof}

The theorem is a limitation of these particular encodings, not a
limitation of every invariant-theoretic construction. A successful
encoding would have to retain constraints absent from row and column
marginals and these constant capacities, such as simultaneous real
sign cancellation between every pair of rows.

\subsection{A quantitative separation despite perfect capacity}
There is a uniform positive separation from flatness in every forbidden
order $n\equiv2\pmod4$, $n>2$, even though the preceding capacities
remain exactly $1$.

\begin{proposition}[Elementary parity gap]
Let $n\equiv2\pmod4$ and $U\in O(n)$. Then
\[
 F_4(U)-1\ge
 \frac{\left(\sqrt{1+\sqrt{2(n-2)/n}}-1\right)^2}{n}
 \ge \frac{2(n-2)}{9n^2}.
\]
Equivalently this lower bounds
$\|U\circ U-J_n/n\|_F^2$, since that squared distance is $F_4(U)-1$.
No optimality of the constants is asserted.
\end{proposition}
\begin{proof}
For any sign matrix $H$ of order $n$, split its rows into two classes
according to the parity of their number of negative entries.
Two rows in the same class have inner product congruent to $2$
modulo $4$, hence of magnitude at least $2$. If one class has $a$
rows, there are $a(a-1)+(n-a)(n-a-1)\ge n(n-2)/2$ ordered pairs of
distinct rows in the same class. Consequently
\[
 \|HH^{\mathsf T}-nI\|_F^2\ge2n(n-2).
\]
Choose $H_{ij}=\operatorname{sign}(u_{ij})$, with either sign allowed
at zero, set $C=H/\sqrt n$, and let $d=\|C-U\|_F$.
Entrywise,
\[
 \left(|u_{ij}|-n^{-1/2}\right)^2
 =\frac{(u_{ij}^2-1/n)^2}{(|u_{ij}|+n^{-1/2})^2}
 \le n(u_{ij}^2-1/n)^2,
\]
so $d^2\le n(F_4(U)-1)$. Writing $C=U+(C-U)$ and using the
triangle inequality and $\|EE^{\mathsf T}\|_F\le\|E\|_F^2$ gives
\[
 L:=\sqrt{2(n-2)/n}\le\|CC^{\mathsf T}-I\|_F\le2d+d^2.
\]
Hence $d\ge\sqrt{1+L}-1$, proving the first bound.
Finally $0\le L<\sqrt2$ and
$(\sqrt{1+L}-1)^2=L^2/(1+\sqrt{1+L})^2\ge L^2/9$.
\end{proof}

For example, at order $6$ every orthogonal matrix passes both capacity
tests with value exactly $1$, yet $F_4(U)\ge1+2/81$.
At admissible orders the same capacity tests stay constant, so they
cannot by themselves distinguish existence from nonexistence of a flat
point. Proving a stronger real-sign-sensitive feasibility theorem remains
an open research task; replacing it by ordinary bistochastic feasibility
would remove the essential difficulty.

\section{A complete hierarchy can still average away the obstruction}
Let $\mathcal H=\mathbb C^n\otimes\mathbb C^n$, $D=n^2$.
For $U\in O(n)$ set $v_U=\operatorname{vec}(U)/\sqrt n$ and
$\sigma_U=v_Uv_U^*$. For any probability measure $\mu$ on $O(n)$,
\[
 \rho_k=\int\sigma_U^{\otimes k}\,d\mu(U)
\]
is permutation invariant and satisfies $\tr_{k+1}\rho_{k+1}=\rho_k$.
Each conditional factor is a pure maximally entangled state, not an
arbitrary relaxed state.
\begin{proposition}[Exact false positive for one-copy uniformity]
For every $n>1$ there is such a hierarchy with $\rho_1=I_D/D$,
yet no factor in its specified representing measure is flat.
\end{proposition}
\begin{proof}
Choose a uniform permutation matrix $P$ and independent uniform row
signs $\epsilon_i\in\{\pm1\}$, and put $U=\operatorname{diag}(\epsilon)P$.
The $(ij),(ab)$ entry of $\mathbb E\sigma_U$ is
$\mathbb E(u_{ij}u_{ab})/n$. Row signs make it zero if $i\ne a$;
if $i=a$ it vanishes unless $j=b$, since a permutation row has only
one nonzero entry. The remaining value is $1/n^2$. Every sampled
$U$ has zeros and so is nonflat. This gives the entire consistent
hierarchy, not merely finitely many symmetric marginals.
\end{proof}
Consequently uniform averaged measurement probabilities, complete
exchangeability, and maximally entangled conditional factors together
do not imply flatness. Indeed $\rho_1$ has maximum entropy $\log D$,
although its displayed ensemble consists solely of sparse matrices.

\section{The precise two-copy obstruction}
Index the standard basis of $\mathcal H$ by $a=(i,j)$ and define
$C=\sum_{a=1}^D |a,a\rangle\langle a,a|$ on $\mathcal H^{\otimes2}$.
For the above hierarchy,
\begin{equation}\label{coincidence}
 \tr(\rho_2 C)=\frac1{n^2}\int F_4(U)\,d\mu(U)
 =\frac1{n^2}\left(1+\int\sum_{ij}(u_{ij}^2-1/n)^2d\mu\right).
\end{equation}
The last identity uses $\sum u_{ij}^2=n$. Thus $\tr(\rho_2C)\ge1/n^2$,
with equality exactly when $\mu$ is supported on flat orthogonals.
For the signed-permutation hierarchy it is instead $1/n$.
This is a nonnegative variance that survives convex averaging.
If the left side is at most $(1+\varepsilon)/n^2$, some factor has
$\sum(u_{ij}^2-1/n)^2\le\varepsilon$; moreover the measure of factors
with defect at least $\delta>0$ is at most $\varepsilon/\delta$.
These statements follow from the integral identity and Markov's
inequality. They provide approximate factors, not exact Hadamards.

\section{An exact quantum moment characterization}
Here an infinitely exchangeable hierarchy means a consistent sequence
of permutation-invariant complex density matrices on
$\mathcal H^{\otimes k}$ for all positive $k$. By the quantum de Finetti
theorem~\cite{CFS}, it has a representation
$\rho_k=\int\sigma^{\otimes k}d\nu(\sigma)$ over density matrices on
$\mathcal H$. Only existence of this representation is used.
We deliberately work over a complex Hilbert space and impose reality
afterwards: the cited paper notes that the corresponding de Finetti
statement fails for real quantum theory, so it cannot be imported
unchanged on a real Hilbert space.
Let $S$ swap the two whole $\mathcal H$ copies, and let $S_A$ swap
only their first $\mathbb C^n$ subsystems. For $a<b$ put
$L_{ab}=i(|a\rangle\langle b|-|b\rangle\langle a|)$, a Hermitian
matrix. All the following constraints concern the same $\rho_2$.
\begin{theorem}[Where the hierarchy becomes equivalent to existence]
A real Hadamard matrix of order $n$ exists if and only if there is
an infinitely exchangeable hierarchy satisfying
\begin{align*}
 \tr(\rho_2S)&=1,& \tr(\rho_2S_A)&=1/n,\\
 \tr(\rho_2(L_{ab}\otimes L_{ab}))&=0\quad(a<b),&
 \tr(\rho_2 C)&=1/n^2.
\end{align*}
Every de Finetti representing measure of such a hierarchy is supported
on vectorized real scaled Hadamard matrices, up to irrelevant phases.
\end{theorem}
\begin{proof}
The swap identity gives $\tr(\rho_2S)=\int\tr(\sigma^2)d\nu$.
Since density-matrix purity is at most one, equality forces $\sigma$
pure almost surely. Similarly
$\tr(\rho_2S_A)=\int\tr(\sigma_A^2)d\nu$. The purity of a density
matrix on $\mathbb C^n$ is at least $1/n$, with equality only at
$I_n/n$, so $\sigma_A=I_n/n$ almost surely.
For each $a<b$ the reality constraint reads
$\int(\tr(\sigma L_{ab}))^2d\nu=0$. The integrand is a square of a
real number, hence vanishes almost surely. The finitely many such
constraints force $\sigma$ real. A real rank-one positive matrix has
a real unit vector representative $v$. Reshape $v$ as an $n\times n$
real matrix $B$. Then $BB^T=I_n/n$, so $U=\sqrt nB\in O(n)$.
Finally $\tr(\rho_2 C)=\int\sum_a\sigma_{aa}^2d\nu$.
Since $\sum_a\sigma_{aa}=1$, this is at least $1/D$ pointwise,
and equality forces all $\sigma_{aa}=1/D$. Thus $u_{ij}^2=1/n$.
The intersection of the finitely many full-measure sets still has
full measure and is nonempty. Conversely a real flat orthogonal
matrix gives the pure tensor hierarchy $\rho_k=\sigma_U^{\otimes k}$,
which satisfies every displayed constraint.
\end{proof}
The theorem is a reformulation and an exact diagnostic, not a proof
of feasibility. In particular the low-order equalities are only used
together with infinite extendibility. We make no claim that an
arbitrary symmetric two-copy positive matrix satisfying the equalities
has a de Finetti representation. Replacing infinite extendibility by
finite symmetry without proof would introduce an additional gap.

\section{Implications for an entropy or tensor-stability bridge}
An optimization over $\rho_1$ alone can reach its ideal entropy in
every dimension using the signed-permutation example; it cannot
force a single flat factor. A bridge must retain a nonlinear
factor-level witness such as the coincidence defect in
\eqref{coincidence}, or another constraint that forces the same
nonnegative variance to vanish. Merely allowing correlations or
symmetric extensions supplies no such vanishing. Even with the exact
two-copy characterization, feasibility at all orders is the original
existence issue in another language.


\section{Which step would actually prove the conjecture?}
The sufficient rounding inequality holds in the requested dimension;
unlike tensor descent it needs no dimension reduction. The missing
step is an existence theorem for distributions meeting that inequality
for every $4\mid n$. Pairwise independence alone is easy, while
concentration near the collision limit is the unresolved constraint.
Optimal min-entropy, correct marginals, positive capacities, modular
orthogonality and one-copy mixed-state entropy each lose information
needed for that concentration or for real sign cancellation.

The completion theorem gives a precise small-defect interface, but
the eight-missing-row example rules out unconditional completion of
arbitrary partial choices beyond it. A promising next question must
therefore select better partial choices, prove a concentration theorem,
or encode additional sign-sensitive constraints in an invariant or
moment formulation. None of these existence statements is proved here.
No Lean formalization, Git commit or publication was made in this pass.

\begin{thebibliography}{9}
\bibitem{GP} D. Gavinsky and P. Pudl\'ak,
On the joint entropy of $d$-wise-independent variables,
Commentationes Mathematicae Universitatis Carolinae 57 (2016), 333--343.
\url{https://users.math.cas.cz/~pudlak/pairwise.pdf}.
\bibitem{Verheiden} E. P. Verheiden,
Arithmetical Properties of Combinatorial Matrices, Ph.D. thesis,
California Institute of Technology (1978), Theorem 4.5, pp. 22--23.
\url{https://thesis.caltech.edu/2899/}.
\bibitem{CFS} C. M. Caves, C. A. Fuchs and R. Schack,
Unknown quantum states: the quantum de Finetti representation,
Journal of Mathematical Physics 43 (2002), 4537--4559.
\url{https://arxiv.org/abs/quant-ph/0104088}.
\end{thebibliography}
\end{document}
